%-------------------------------------------------------------------------------
\section{Payment Protocol Modeling}
%-------------------------------------------------------------------------------
We formally model the payment protocols of WeChat Pay and Alipay across three payment scenarios, and the complete models are available in an open-source repository~\cite{Repo2025FormalModels}.
In this section, we present the key details in modeling third-party payment protocols.
First, we introduce how a generalized channel model is incorporated into third-party payment protocols to flexibly capture different channel assumptions.
Then, we describe the user model, which allows us to accurately capture users' evaluation of payment information and their authorization behavior within the symbolic model.
Finally, we explain how to use an Order Oracle to model the adversary’s capability to obtain legitimate payment parameters.

\subsection{Channel Model}
We introduce appropriate modifications to the standard secure channel model in order to systematically model different channels and security assumptions.
We use the facts \TFact{SOut} and \TFact{SIn} to represent message sending and receiving over a channel.
Messages transmitted over our secure channels are not replayable.
Therefore, the message transmission process of the channel is simplified, and the rule is shown below.
\begin{flalign*}
\text{\Tterm{Ru}}  & \text{\Tterm{le SecureChannel}} :      & \\
  \TFacts{  & \text{\TFactTerm{SOut}{ChannelID, A, B, m}}}  \\
            & \TEmpActFactSL{}                                \\
  \TFacts{  & \text{\TFactTerm{SIn}{ChannelID, A, B, m}}}
\end{flalign*}
Here, \Tterm{m} is the message transmitted over the channel, \Tterm{A} denotes the message sender, and \Tterm{B} denotes the message receiver.
\Tterm{ChannelID} is a channel identifier that specifies which channel is used for communication.
For example, when \Tterm{ChannelID} takes the value \Tterm{\textquotesingle U-PS\textquotesingle}, it denotes the communication channel between \emph{U} and \emph{PS}.

To model insecure channels, we introduce two additional attack rules for such channels, corresponding to violations of channel confidentiality and integrity, respectively.
\begin{flalign*}
\text{\Tterm{Ru}}  & \text{\Tterm{le BreakConfidentiality}} :                   & \\
  \TFacts{  & \text{\TFactTerm{SOut}{ChannelID, A, B, m}} }                     \\
            & \TActFactsSL{ \text{\TFactTerm{InsecureChannel}{ChannelID}}}      \\
  \TFacts{  & \text{\TFactTerm{SOut}{ChannelID, A, B, m}, \TFactTerm{Out}{m}}}  \\ \\
\text{\Tterm{Ru}}  & \text{\Tterm{le BreakIntegrity}} :                         \\
  \TFacts{  & \text{\TFactTerm{In}{\TPairN{ChannelID, A, B, m}}} }              \\
            & \TActFactsSL{ \text{\TFactTerm{InsecureChannel}{ChannelID}}}      \\
  \TFacts{  & \text{\TFactTerm{SOut}{ChannelID, A, B, m}}}                      
\end{flalign*}
Compromising confidentiality grants the adversary the ability to obtain messages from the channel, while compromising integrity grants the adversary the ability to inject arbitrary messages into the channel.
\TFact{In} and \TFact{Out} are built-in facts of the Tamarin prover that represent receiving messages from, and sending messages to, the public channel.
The fact \TFact{InsecureChannel} is used to label a channel as insecure. 
We use a restriction to constrain the adversary so that it can compromise only the channels specified under a given channel assumption.
\begin{flalign*}
\text{\Tterm{Re}}  & \text{\Tterm{striction Insecure}} :            & \\
    & \TAll \; \text{\Tterm{ch\_id k \#i.}} \; \text{\TFact{InsecureChannel(ch\_id)@\#i}} \\
    & \implies (\text{\Tterm{ch\_id}} = \text{\Tterm{\textquotesingle MF-MS\textquotesingle}} \; \lor \; \text{\Tterm{ch\_id}} = \text{\Tterm{\textquotesingle MF-PS\textquotesingle}} ) 
\end{flalign*}
This restriction specifies that the channels between \emph{MF} and \emph{MS}, and between \emph{MF} and \emph{PS}, are insecure.

\subsection{User Model} 
Under our security assumptions, users are capable of accurately identifying the payment amount and product description, thereby avoiding being misled into authorizing a payment.
To model this capability in the symbolic model, we introduce two unary functions, 
\Tfunction{calc\_amt} and \Tfunction{gen\_desc}, which are used to compute the payment amount and the product description from goods information, respectively.
During the \emph{Ordering} phase, users select goods.
In our model, we use a fresh value as an abstract representation of the goods information.
The following rule models the user's goods ordering behavior:
\begin{flalign*}
      &\TFacts{ \text{\TFact{Fr}}(\TSymbolFr goods)} \TEmpActFactSL{} & \\
      &\TFacts{
        \text{\TFact{Order}}(\TSymbolFr goods),
        \text{\TFact{SOut}}(..., \TSymbolFr goods)
      }
\end{flalign*}
This rule represents that the user \emph{U} selects a product $\TSymbolFr goods$ and sends it to \emph{MF}.
Here, $\text{\TFact{Fr}}(\TSymbolFr goods)$ indicates that the value $\TSymbolFr goods$ is freshly generated.
% 我们简化了通信
In \emph{Prepayment} phase, \emph{MS} computes the payment $amount$ and $description$ based on $\TSymbolFr goods$ using \Tfunction{calc\_amt} and \Tfunction{gen\_desc}, as follows:
\begin{flalign*}
      amount & = \text{\Tfunction{calc\_amt}}(\TSymbolFr goods) &\\
      description & = \text{\Tfunction{gen\_desc}}(\TSymbolFr goods)
\end{flalign*}
In \emph{Payment} phase, the user obtains the order information from the \emph{PS-App} and confirms it.
Due to pattern matching, the user authorizes a payment if only if the payment $amount$ and $description$ exactly match $\text{\Tfunction{calc\_amt}}(\TSymbolFr goods)$ and $\text{\Tfunction{gen\_desc}}(\TSymbolFr goods)$, respectively, and will not authorize any other values.
The following rule models the user's payment authorization behavior:
\begin{flalign*}
      \TFacts{ &\text{\TFact{Order}}(\TSymbolFr goods),                     &\\
              &\text{\TFact{SIn}}(..., \TPair{\text{\Tfunction{calc\_amt}}(\TSymbolFr goods), \text{\Tfunction{gen\_desc}}(\TSymbolFr goods)}) } \\
      & \TEmpActFactSL{} \\
      \TFacts{& \text{\TFact{UserAuthorizePayment}}(...) }
\end{flalign*}

\subsection{Order Oracle Modeling}

\begin{figure}[t]
  \centering
  \includegraphics[width=\linewidth]{img/black.png}
  \caption{Overview of the Order Oracle}
  \label{fig:order_oracle_overview}
\end{figure}

In our threat model, the adversary is assumed to have the capability to obtain legitimate payment parameters from arbitrary merchant platforms.
To model this capability in an intuitive and explicit manner, we introduce the concept of an \emph{Order Oracle}.
Figure~\ref{fig:order_oracle_overview} illustrates the structure of the \emph{Order Oracle}.
\emph{Order Oracle} takes two inputs: the payment amount $amount$ and the order description $description$.
Upon receiving these two inputs, \emph{Order Oracle} generates payment parameters $pay\_params$ and returns them to the adversary.
The generated $pay\_params$ constitute valid payment parameters accepted by the \emph{PS}.
When $pay\_params$ are used for payment, the payment amount and order description displayed on the payment interface are consistent with the provided $amount$ and $description$.
We model the Order Oracle using the following rule:
\begin{flalign*}
\text{\Tterm{Ru}}  & \text{\Tterm{le OrderOracle}} :        & \\
  \TFacts{  & \text{\TFactTerm{In}{<amount, description>}}} \\
            & \TEmpActFactSL{}                              \\
  \TFacts{  & \text{\TFactTerm{Out}{\text{\Tfunction{gen\_params(amount, description)}}}}}
\end{flalign*}
\Tfunction{gen\_params} is a function that generates valid payment parameters $pay\_params$ based on the payment amount $amount$ and the order description $description$.
This function is not deterministic. 
In our formal model, \Tfunction{gen\_params} represents different payment-parameter generation procedures under different payment scenarios and protocols, 
thereby producing payment parameters that are compliant with the corresponding protocol specifications.